\documentclass[compacto,10pt]{pubciencias-notas}

% -- Paquetes adicionales
\usepackage{comandosEPN}

% -- Datos del libro
\autor{Departamento de Formación Básica}
\materia{Álgebra Lineal}
\nota{Semana 02 Clase 02}
\fecha{Semestre 2025-A}

\logouno[1.1cm]{Logo_EPN.eps}
\logodos[1.7cm]{Logo_DFB.eps}

\definecolor{colordef}{cmyk}{0.81,0.62,0.00,0.22}

\begin{document}

\encabezado

\vspace{-8mm}

\section{Matrices}
\setcounter{subsection}{3}
\subsection{Equivalencia por filas}

\textbf{Pasos para transformar una matriz a una matriz escalonada por filas.} Dada una matriz $A$.
\begin{enumerate}
\item Determinar la primera columna (contando de izquierda a derecha), tal que no todos sus elementos sean cero. Ésta es la \textbf{columna pivote}.
\item En la columna pivote, identificar al primer elemento (contando de arriba hacia abajo) distinto de cero. A este elemento lo denominaremos \textbf{pivote}.
\item Intercambiar, en caso necesario, la primera fila por aquella fila donde aparece el pivote, de modo que éste se encuentre ahora en la primera fila.
% \item Multiplicar a la primera fila por el recíproco del pivote. Así, el primer elemento de la columna pivote es ahora un 1.
\item Realizar las operaciones elementales de fila necesarias para que los elementos debajo del pivote, dentro de la columna pivote, sean cero.
\item Identificar a la primera columna (después de la columna pivote) cuyos elementos sean no todos cero y repetir los pasos anteriores hasta obtener a la matriz escalonada por filas.
\end{enumerate}

\begin{ejem}\hspace{0pt}\label{Ejem:EPF} Dada la matriz
\[
    A=\begin{pmatrix}
    1&5&-4\\
    1&-2&1\\
    3&1&-2
    \end{pmatrix}.
\]
Obtengamos una matriz escalonada por filas, para esto, notemos que la columna pivote es la primera columna de $A$ y que en la posición $11$ ya se encuentra el pivote, entonces
\begin{align*}
A=\begin{pmatrix}
1&5&-4\\
1&-2&1\\
3&1&-2
\end{pmatrix}
& \sim 
\begin{pmatrix}
1&5&-4\\
0&-7&5\\
0&-14&10
\end{pmatrix}
&& \begin{array}{c}
    -F_1+F_2\rightarrow F_2\\
    -3F_1+F_3\rightarrow F_3
\end{array}\\
& \sim 
\begin{pmatrix}
1&5&-4\\
0&-7&5\\
0&0&0
\end{pmatrix}
&& -2F_2+F_3\rightarrow F_3
% & \sim 
% \begin{pmatrix}
% 1&5&-4\\
% 0&1&-\frac{5}{7}\\
% 0&0&0
% \end{pmatrix}
% && -\frac{1}{7}F_2\rightarrow F_2
\end{align*}
Luego, una matriz escalonada por filas obtenida a partir de $A$ es
\[
    \begin{pmatrix}
    1&5&-4\\
    0&-7&5\\
    0&0&0
    \end{pmatrix}.
\]
\end{ejem}

\begin{ejer}\hspace{0pt} Dada la matriz
\begin{align*}
A=\begin{pmatrix}
1&5&-4&0\\
1&-2&1&2\\
0&1&-1&1
\end{pmatrix}.
\end{align*}
Obtenga una matriz escalonada por filas equivalente.
\end{ejer}

\begin{teo}
    Toda matriz $A\in\Mat[\K]{m}{n}$ es equivalente por filas a una única matriz en forma escalonada reducida por filas. 
\end{teo}

\textbf{Pasos para transformar una matriz a una matriz escalonada reducida por filas.} Dada una matriz $A$.
\begin{enumerate}
\item Determinar la primera columna (contando de izquierda a derecha), tal que no todos sus elementos sean cero. Ésta es la \textbf{columna pivote}.
\item En la columna pivote, identificar al primer elemento (contando de arriba hacia abajo) distinto de cero. A este elemento lo denominaremos \textbf{pivote}.
\item Intercambiar, en caso necesario, la primera fila por aquella fila donde aparece el pivote, de modo que éste se encuentre ahora en la primera fila.
\item Multiplicar a la primera fila por el recíproco del pivote. Así, el primer elemento de la columna pivote es ahora un 1.
\item Realizar las operaciones elementales de fila necesarias para que los elementos de la columna pivote sean cero, excepto el pivote.
\item Identificar a la primera columna (después de la columna pivote) cuyos elementos sean no todos cero y repetir los pasos anteriores de tal manera que ahora solamente, en la nueva columna pivote, su elemento pivote sea diferente de cero. Continuar con el proceso indicado hasta obtener a la matriz escalonada reducida por filas.
\end{enumerate}

%Hallar una matriz $B$ equivalente a la matriz $A$ mediante el proceso de Gauss.

%\begin{advertencia}
%    \textbf{Eliminación de Gauss Jordán:} El proceso para obtener una matriz escalonada reducida por filas a partir de una matriz cualquiera se conoce como eliminación de Gauss-Jordán.
%\end{advertencia}

\begin{ejem}\label{Ejem:EPF}\hspace{0pt} Dada la matriz
\[
    A=\begin{pmatrix}
    1&5&-4\\
    1&-2&1\\
    3&1&-2
    \end{pmatrix}.
\]
Obtengamos la correspondiente matriz escalonada reducida por filas equivalente. Se tiene que
\begin{align*}
A=\begin{pmatrix}
1&5&-4\\
1&-2&1\\
3&1&-2
\end{pmatrix}
& \sim 
\begin{pmatrix}
1&5&-4\\
0&-7&5\\
0&-14&10
\end{pmatrix}
&& \begin{array}{c}
    -F_1+F_2\rightarrow F_2\\
    -3F_1+F_3\rightarrow F_3
\end{array}\\
& \sim 
\begin{pmatrix}
1&5&-4\\
0&-7&5\\
0&0&0
\end{pmatrix}
&& -2F_2+F_3\rightarrow F_3\\
& \sim 
\begin{pmatrix}
1&5&-4\\
0&1&-\frac{5}{7}\\
0&0&0
\end{pmatrix}
&& -\frac{1}{7}F_2\rightarrow F_2\\
& \sim
\begin{pmatrix}
1&0&-\frac{3}{7}\\
0&1&-\frac{5}{7}\\
0&0&0\\
\end{pmatrix}
&&-5F_2+F_1\rightarrow F_1
\end{align*}
Luego, la correspondiente matriz escalonada por filas obtenida a partir de $A$ es
\[
    \begin{pmatrix}
    1&0&-\frac{3}{7}\\
    0&1&-\frac{5}{7}\\
    0&0&0\\
    \end{pmatrix}.
\]
\end{ejem}

\begin{ejer}
Dada la matriz:
\[
    A=\begin{pmatrix}
    1&5&-4&0\\
    1&-2&1&2\\
    0&1&-1&1
    \end{pmatrix}
\]
Obtenga su correspondiente matriz escalonada reducida por filas.
\end{ejer}

\begin{defi}
    Sean $A, B\in\Mat[\K]{m}{n}$, donde 
    %%%%%%%%%%%%%%%%%%%%%%%%%%%%%%%%%%%%%%%%%%%%%%%%%%%%%%%%%%%%%%%%
    %la única matriz escalonada reducida por filas 
    %%%%%%%%%%%%%%%%%%%%%%%%%%%%%%%%%%%%%%%%%%%%%%%%%%%%%%%%%%%%%%%%
    $B$ es una matriz escalonada por filas 
    equivalente a $A$. 
    El rango de $A$, denotado por $\rang(A)$, es el número de filas no nulas que tiene la matriz $B$.  
\end{defi}

\begin{prop}
    Sean $A, B\in\Mat[\K]{m}{n}$. Se tiene que si $A\sim B$, entonces 
    \[
        \rang(A) = \rang(B).
    \]
\end{prop}

\begin{prop}
    Sea $A\in\Mat[\K]{m}{n}$. Se tiene que si $A$ es una matriz escalonada, entonces $\rang(A)$ es el número de filas no nulas que tiene $A$.
\end{prop}

\begin{ejem}\hspace{0pt} Considerando que
\[
    A=\begin{pmatrix}
        1&5&-4\\
        1&-2&1\\
        3&1&-1
    \end{pmatrix} \sim
    \begin{pmatrix}
        1&5&-4\\
        0&-7&5\\
        0&0&0
    \end{pmatrix}=B
\]
y que $B$ es una matriz escalonada con dos filas no nulas, entonces $\rang(A)=\rang(B)=2$
\end{ejem}

\begin{ejer} \label{Ejer:OpeELe} Calcular el rango de la matriz 
\[
    B=\begin{pmatrix}
    1&0&-2&0\\
    1&-2&1&-1\\
    -2&1&3&4
    \end{pmatrix}.
\]
\end{ejer}
\begin{defi}[Matriz elemental]
    Dada una operación elemental por fila, se llama matriz elemental al resultado de aplicar dicha operación a la matriz identidad.
\end{defi}

\begin{ejem}\label{ejem2}
Considerando la matriz identidad $I_{3}$ y la operación elemental de fila $F_3+F_2\rightarrow F_2$, tenemos que le corresponde la matriz elemental:
\[
    E=\begin{pmatrix}1&0&0\\0&1&1\\0&0&1\end{pmatrix}.
\]
\end{ejem}

\begin{ejer}\hspace{0pt} Considere las siguientes operaciones elementales de fila y escriba las correspondiente matrices elementales.
\begin{enumerate}
    \item $-F_1+F_2\rightarrow F_2$
    \item $2F_1+F_3\rightarrow F_3$
    \item $F_2\longleftrightarrow F_1$.
\end{enumerate}
\end{ejer}

\begin{teo}
    Sean $E\in \Mat[\K]{m}{m}$ una matriz elemental y $A\in\Mat[\K]{m}{n}$, el resultado de aplicar la operación por filas correspondiente a $E$ a la matriz $A$ es $EA$.
\end{teo}

\begin{ejem}\hspace{0pt} Considere a la matriz 
\[
    A=\begin{pmatrix}0&2&5\\-4&-1&0\\3&2&1\end{pmatrix}
\]
y a la operación elemental de fila $F_3+F_2\rightarrow F_2$. Notemos que al aplicar esta operación a $A$ se obtiene la matriz equivalente $B$.
\begin{align*}
	A=\begin{pmatrix}
	0&2&5\\
	-4&-1&0\\
	3&2&1\\
	\end{pmatrix}
	\sim
	 \begin{pmatrix}
    0&2&5\\
    -1&1&1\\
    3&2&1\\
    \end{pmatrix} = B
\end{align*}
Por otro lado, se tiene que la matriz elemental correspondiente a dicha operación es
\[
    E=
    \begin{pmatrix}1&0&0\\0&1&1\\0&0&1\end{pmatrix}
\]
tal que 
\[
    B=EA=\begin{pmatrix}1&0&0\\0&1&1\\0&0&1\end{pmatrix}\begin{pmatrix}
	0&2&5\\
	-4&-1&0\\
	3&2&1\\
	\end{pmatrix}=\begin{pmatrix}
    0&2&5\\
    -1&1&1\\
    3&2&1\\
    \end{pmatrix}
\]
\end{ejem}

\begin{ejer}\hspace{0pt} Dada la matriz $A = \begin{pmatrix}2&-3\\2&1\end{pmatrix}$ considere a las matrices elementales
\[
    E_{1} = \begin{pmatrix}1&0\\-1&1\end{pmatrix}
    \texty
    E_{2} = \begin{pmatrix}\frac{1}{2}&0\\0&1\end{pmatrix}
\]
tal que la matriz $B = E_{2}E_{1}A$.
\begin{enumerate}
\item[$a)$] Encuentre a la matriz $B$.
\item[$b)$] Identifique las operaciones elementales de fila asociadas a las matrices elementales $E_{1}$ y $E_{2}$.
\item[$c)$] Aplique las operaciones elementales de fila determinadas en el literal anterior sobre A para mostrar que $A\sim B$.
\end{enumerate}
\end{ejer}

\end{document}
